\documentclass[11pt]{article}
\usepackage[T1]{fontenc}
\usepackage[utf8]{inputenc}
\usepackage{lmodern,amsmath,amssymb,amsthm,microtype,booktabs,xcolor}
\usepackage[margin=1in]{geometry}
\usepackage[colorlinks=true,linkcolor=blue!50!black,urlcolor=blue!50!black]{hyperref}
\setlength{\emergencystretch}{3em}
\newtheorem{proposition}{Proposition}
\title{Proof audit and scope of the scaling results}
\author{ChatGPT-6 Astra}
\date{October 15, 2026}
\begin{document}
\maketitle

This report records the findings from checking the supplied Fable draft and its cited mean-field argument. The accompanying papers give the replacement proofs in full. All learning statements below use the prior conditioned on an absolute output margin at \emph{every} Boolean input, followed by noiseless observed-label conditioning. A confidence-success statement means that the posterior probability of the correct label is strictly greater than $2/3$.

\section{Audit of the supplied finite-exponent manuscript}
\label{sec:fable-audit}

We checked the supplied manuscript by Claude Fable 5, \emph{The circuit
prior across the whole strong prior regime: cloning and subsampling at
every $\gamma>1$}. We found no error invalidating its central finite-$\gamma$ proof. In
particular, the Gaussian change of measure, robustness of the cloned
witness, singular-vector frame net, orthonormalization of the net points,
and spectral truncation preserve the stated margin and cost bounds.
The conversion from a low-cost representative to a small bounded-weight
network uses the compression theorem proved in the earlier paper.

The main improvement concerns the learning theorem. The supplied proof
bounds the tail by multiplying the largest excluded function probability
by $2^{2^d}$, the number of Boolean functions. This gives a valid but
unnecessary additional requirement $N^\gamma\gtrsim2^d$. Its own
structural estimate yields a stronger conclusion: outside one exceptional
parameter event, \emph{every} globally margin-separated network has a
low-cost representative. The total probability of functions without such
a representative is therefore bounded by the exceptional parameter
probability directly. Counting only the resulting low-cost class gives
the width-uniform learning theorem in the accompanying paper. The source's
exponential width requirement is thus a limitation of its tail estimate,
not a demonstrated obstruction to learning.

We also replace the singular-vector net with Gaussian enlargement
applied to distance from an operator-norm ball. This shortens the proof
and removes its rank and net-entropy bookkeeping. A direct chi-square
estimate removes the logarithmic width proviso in its small-width
complexity comparison.

Two qualifications concern the source's discussion rather than its main
proof. First, the simulation in which a spectral cut lies below the noise
bulk demonstrates failure of that particular cut. It does not establish
necessity of the theorem's quantitative width threshold, nor a lower
bound on the width required for generalization. Second, the conjectured
$\gamma=1$ extension is not proved by the source's vanishing-operator-noise
argument. Any endpoint result needs a separate argument; the endpoint
statements in the accompanying endpoint paper are proved independently.

Minor points do not affect the results: ceiling operations in the
small-width bound require an additive one or an absolute factor; the
rank estimate is $r_*\le18K/\delta^2$; and the no-inner-layer case $L=2$
is most clearly handled as an empty truncation sum. We found no major
mistake invalidating the supplied manuscript's finite-$\gamma>1$
complexity theorem.


\section{The cited mean-field proof: a definite pointwise-norm error}
The finite-exponent draft cites Vaintrob and H\"anni, \emph{Width-Robust Learnability in Mean-Field Bayesian Neural Networks}, \href{https://arxiv.org/abs/2607.05735}{arXiv:2607.05735v1}. Its model uses a smooth squared-error Gibbs tilt and independent biases at every hidden layer. Those conventions differ from the globally margin-conditioned classifier considered here, so its learning conclusion could not be transferred without additional work even if its proof were complete.

There is also a direct mathematical error in its proposed regularity invariant. Appendix A, Definition A.6(I4), bounds a pointwise derivative by a population $L^2$ norm of the perturbation, with a polynomial constant. The claimed initialization in Appendix B, Theorem B.14, fails already for a simple readout.

Let $\mu$ be uniform on $\{-1,1\}^d$ and consider
\[
 S(\xi)(x)=\frac1M\sum_{j=1}^M\tanh(\xi_j(x)).
\]
The asserted form of the bound is
\[
 |D_{j,u}S(\xi)(x)|\leq\frac A{\sqrt M}\|u\|_{L^2(\mu)}
 \quad\text{for every direction }u.
\]
At $\xi=0$, choose $u=2^{d/2}{\bf1}_{\{x\}}$. Then
\[
 \|u\|_{L^2(\mu)}=1,\qquad
 D_{j,u}S(0)(x)=\frac{2^{d/2}}M.
\]
The estimate forces $A\geq2^{d/2}/\sqrt M$. Thus polynomial width and a polynomial regularity constant cannot satisfy the stated invariant. Replacing the population norm by a pointwise norm, or replacing pointwise output control by an averaged loss bound, changes the induction that would need to be proved.

This is a defect in the proof as written. It is not a counterexample to the desired mean-field learning theorem. The new endpoint paper establishes an independent theorem for tanh under the requested posterior convention, using an exact finite-domain variational principle and compression of its minimizers.

\section{Two further steps requiring repair in that appendix}
\subsection{A second-moment tail does not follow from bounded entropy}
Proposition A.7 uses bounded relative entropy and second moments to conclude that the second moment restricted to a sufficiently large exceptional set can be made uniformly small. The stated inference does not hold in that generality.

For $b>1$, put
\[
 Q_b=(1-b^{-2})\mathcal N(0,1)+b^{-2}\mathcal N(b,1).
\]
Convexity of relative entropy gives
\[
 D(Q_b\Vert\mathcal N(0,1))\leq
 b^{-2}D(\mathcal N(b,1)\Vert\mathcal N(0,1))=\frac12.
\]
The second moments are uniformly bounded, but
\[
 \liminf_{b\to\infty}
 \mathbb E_{Q_b}[X^2{\bf1}_{\{|X|>b/2\}}]\geq1.
\]
Thus neither bounded relative entropy nor the second-moment estimate used there supplies a uniform-integrability rate for the second moment. This diagnoses the displayed justification; it does not rule out a different truncation argument using the bounded activation or additional properties of the particular Gibbs law.

\subsection{Fresh Gaussian noise must be accounted for}
The explanation accompanying Lemma B.13 treats a realized reconstructed network as if replacing its fresh Gaussian contribution by a conditional expectation introduced no additional fluctuation. Independence of the fresh noise does not provide this conclusion. For a simple test, take fixed nonzero $a$ and independent $\xi_i\sim\mathcal N(0,1)$. Then
\[
 V_N=\frac a{\sqrt N}\sum_{i=1}^N\xi_i
 \ \stackrel{\mathrm{law}}=\ aZ,\qquad Z\sim\mathcal N(0,1).
\]
The random output $\tanh(V_N)$ has a nonvanishing variance for every $N$. Increasing the number of summands improves its covariance approximation but does not make that realized field concentrate about its mean. If a suffix averages $M$ independent such rows, its fluctuation is governed by that averaging as well, not solely by the lower-layer width $N$.

Accordingly, a proof passing from a comparison of expected suffix outputs to a pointwise finite-network approximation must track the fresh-noise fluctuations and the dependence of their covariance on the sampled features. The observation above isolates the missing step in the stated reasoning; it is not asserted as a counterexample satisfying every auxiliary regularity condition of the final theorem. The new compression proof avoids this issue by concentrating empirical feature covariance matrices and evaluating a bounded readout directly.

\section{What the replacement proofs establish}
\paragraph{Every fixed $\gamma>1$, arbitrary fixed depth.}
For any bounded $1$-Lipschitz activation and a target with finite bounded-weight network complexity $C$, the new strong-scaling paper proves polynomial sample complexity uniformly for every learner width $N\geq C$. The polynomial degree may depend on $\gamma$ and the depth. A common cylinder construction supplies the target-mass bound for every $\gamma\geq1$.

\paragraph{The mean-field endpoint $\gamma=1$.}
For tanh at arbitrary fixed depth, the endpoint paper derives the exact speed-$N$ rate by Sanov's theorem and continuous contraction. A minimum-rate full-margin representation has feature laws tilted by bounded quadratic forms. Output optimality bounds the readout coefficients in $\ell^1$; Gaussian integration by parts controls all earlier quadratic forms. A dimension-independent covariance estimate then gives polynomial-width, bounded-weight realizations of all bounded-rate Boolean functions. Counting these functions proves polynomial sample complexity for every sufficiently large width. The proof does not give a quantitative bound on that width threshold. For one hidden layer, a separate proof gives the stronger every-width result for all $\gamma\geq1$, under the general bounded Lipschitz activation assumption.

\paragraph{Below the critical exponent.}
For tanh and each fixed $\gamma<1$, the hidden kernel freezes to an explicit iterated-arcsine kernel. The posterior limit is a minimum-full-margin-energy classifier for $0<\gamma<1$, a globally margin-conditioned Gaussian classifier for $\gamma=0$, and a Gaussian sign classifier for $\gamma<0$. For $\gamma\leq0$ at a specified depth $O(d)$, the new confidence criterion gives an exact expected failure rate $(1-2^{-d})^n$; taking $d=k$ yields exponential sample complexity for parity. The universal claim for every fixed depth and every $\gamma<1$ remains unresolved here. In particular, the $0<\gamma<1$ classification problem is not settled by the regression-mean parity eigenvalue argument.

\section{Effect of the confidence convention}
Let $R$ be average posterior wrong-label probability and let $F$ be the fraction of test inputs where correct-label probability is at most $2/3$. For every posterior,
\[
 \max\{0,(3R-1)/2\}\leq F\leq\min\{1,3R\}.
\]
The positive learning results therefore imply the requested confidence guarantee with a constant-factor change in accuracy. The near-independent NNGP example gives more: every unseen input fails the confidence test, and every seen input passes, so $F$ is exactly the unseen fraction. This resolves the possible tiny-majority-bias ambiguity in that example. It does not convert all small-eigenvalue statements into classification sample lower bounds.

The independent checks covered the cylinder construction, spectral tail, aggregate probability bound, finite-class counting, critical variational representation, minimizer stationarity, covariance comparison, forward compression, and the subcritical posterior limits. Mathematical conclusions in the papers are distinguished from remaining questions; compilation and numerical sanity checks were used only to check presentation and catch inconsistencies.
\end{document}
